\section{Обзор существующих решений}
\label{sec:Chapter3} \index{Chapter3}

Метод SLP получил широкое распространение в современных оптимизирующих компиляторах. Наиболее известные промышленные реализации присутствуют в инфраструктурах LLVM и GCC. За время развития данных проектов алгоритмы SLP-векторизации существенно усложнились и были дополнены многочисленными эвристиками, ориентированными на повышение качества генерируемого кода и поддержку широкого спектра архитектур.

\subsection{Реализация в LLVM}

SLP-векторизатор был добавлен в LLVM в 2013 году. Первоначальная реализация представляла собой сравнительно компактный алгоритм, ориентированный на поиск групп изоморфных инструкций и построение деревьев векторизации внутри базовых блоков.

Основными возможностями ранней реализации являлись:
\begin{itemize}
    \item объединение однотипных арифметических операций;
    \item векторизация последовательных загрузок и сохранений;
    \item базовая модель оценки прибыльности преобразований.
\end{itemize}

Дальнейшее развитие LLVM привело к существенному расширению функциональности SLP-векторизатора. Современная реализация поддерживает:
\begin{itemize}
    \item более сложные схемы перестановки элементов;
    \item частичную векторизацию;
    \item сложные стратегии переупорядочивания операций.
\end{itemize}

В результате кодовая база современного SLP-векторизатора LLVM значительно выросла и включает большое количество эвристик, ориентированных на конкретные архитектурные особенности.

В рамках данной работы для анализа была рассмотрена ранняя реализация SLP-векторизации, представленная в LLVM версии 3.3. Выбор именно данной версии обусловлен несколькими причинами:
\begin{enumerate}
    \item Ранняя реализация содержит базовые алгоритмические идеи метода в относительно компактной и понятной форме.
    \item Она лишена значительного количества инженерных усложнений, накопленных в современных версиях LLVM.
    \item Архитектурные решения раннего SLP-векторизатора лучше подходят для адаптации в условиях ограниченной инфраструктуры анализа, характерной для компилятора Go.
\end{enumerate}

Таким образом, использование реализации LLVM 3.3 для анализа позволяет опираться на проверенный промышленный подход, сохраняя при этом приемлемую сложность.

\subsubsection{Архитектура раннего SLP-векторизатора LLVM}

В основе алгоритма лежит обход базового блока, начинающийся с поиска цепочек последовательных инструкций сохранения. Для этого используется анализатор скалярной эволюции (ScalarEvolution~\cite{ScalarEvolution}), определяющий, обращаются ли две инструкции к смежным областям памяти. Если две записи сохраняют однотипные значения по адресам, отличающимся ровно на размер элемента, они считаются последовательными.

Найденные пары объединяются в цепочки: пока следующее сохранение является смежным с предыдущим и ещё не помечено как векторизованное, оно добавляется в группу. Далее для каждой цепочки запускается попытка векторизации скользящим окном, размер которого определяется целевой шириной векторного регистра (по умолчанию 128 бит).

\subsubsection{Построение и оценка дерева выражений}

Для выбранного окна сохранений строится дерево выражений путём рекурсивного подъёма по аргументам. Алгоритм поднимается к операциям, породившим сохраняемые значения. Рекурсивно формируется набор инструкций-предшественников, проверяется их изоморфизм и допустимость упаковки.

Процесс управляется классом \texttt{BoUpSLP}, который на этапе анализа стоимости просматривает дерево с ограничением глубины. Все инструкции нумеруются для контроля порядка, а также фиксируется, в какой «линии» (позиции внутри будущего вектора) используется каждое значение. Если одна и та же инструкция оказывается востребованной в разных линиях, она помечается как подлежащая обязательной скаляризации~--- её результат не может быть частью векторного регистра без дублирования, что обычно дорого.

Стоимость векторизации рассчитывается рекурсивно: для каждого узла дерева стоимость скалярных эквивалентов сравнивается со стоимостью векторной операции, получаемой через интерфейс \texttt{TargetTransformInfo}. Учитываются затраты на загрузку/сохранение, арифметические операции и возможные вставки/извлечения элементов. Дополнительно проверяется, не мешают ли промежуточные инструкции, работающие с памятью, безопасному перемещению ранних нагрузок к месту размещения векторизованных операций (анализ \texttt{isUnsafeToSink}).

\subsubsection{Генерация векторного кода}

Если модель стоимости показывает выгоду, запускается \texttt{vectorizeTree}. Рекурсивно, начиная с листьев дерева (обычно загрузок или констант) и двигаясь к корню, создаются векторные инструкции.

Для случаев, когда дерево не завершается сохранением, существует метод \texttt{vectorizeArith}, который строит вектор, а затем вставляет операции извлечения элементов для каждого исходного использования, сохраняя корректность программы.

\subsubsection{Основные ограничения ранней реализации}

Ранний SLP-векторизатор LLVM обладает рядом естественных ограничений, многие из которых были сняты в более поздних версиях:
\begin{itemize}
    \item \textbf{Один базовый блок}: векторизация не переходит через границы блоков; PHI-узлы не поддерживаются.
    \item \textbf{Фиксированная векторная ширина}: используется минимальный размер векторного регистра (128 бит), кратность ширине скалярного типа. Нет поддержки частичной упаковки неполных векторов.
    \item \textbf{Жёсткий анализ последовательности}: функция \texttt{isConsecutiveAccess} работает только с константным расстоянием между адресами, не поддерживая индексацию с переменным шагом.
    \item \textbf{Глубина рекурсии}: ограничение в несколько уровней рекурсии не позволяет векторизовать глубокие деревья выражений.
    \item \textbf{Правило однократного использования}: если скалярная инструкция имеет несколько использований, она помечается \texttt{MustScalarize}, что часто делает векторизацию невыгодной.
\end{itemize}

\subsubsection{Значение для выполненной работы}

Ранняя реализация LLVM примечательна тем, что в сжатой форме воплощает ключевые принципы SLP: поиск начальных пар через память, рекурсивное расширение, оценку стоимости и генерацию векторных инструкций. Её изучение позволяет выделить алгоритмическое ядро, свободное от инженерных наслоений поздних версий. Именно поэтому версия LLVM 3.3 выбрана в качестве основного референса для проектирования собственного прохода.

Вместе с тем, восходящий от сохранений подход плохо согласуется с устройством SSA-формы компилятора Go, где целостность памяти задана явным потоком значений \texttt{Memory}. Прямое следование логике LLVM потребовало бы существенной переработки, что выводится за рамки данного этапа и подробно обсуждается в разделе~\ref{sec:go_limitations} (Ограничения переноса существующих решений в компилятор Go).

\subsection{Реализация в GCC}

В компиляторе GCC SLP-векторизация реализована как часть более общей инфраструктуры loop и basic-block vectorization. Подход GCC тесно интегрирован с развитой системой анализа циклов и зависимостей. Это позволяет эффективно сочетать цикловую и SLP-векторизацию, однако требует значительно более сложной инфраструктуры промежуточного представления и анализа.

\emph{Детальный анализ реализации GCC выходит за рамки данной работы, так как основным референсом выбран LLVM. При необходимости описание может быть добавлено на следующем этапе.}

\subsection{Ограничения переноса существующих решений в компилятор Go}
\label{sec:go_limitations}

Несмотря на наличие работоспособных и проверенных временем реализаций SLP-векторизации в таких компиляторах, как LLVM и GCC, их прямое или близкое к исходному заимствование в компиляторе Go оказывается невозможным. Причины лежат в фундаментальных различиях архитектуры промежуточного представления, состава аналитической инфраструктуры и целевых установок проекта.

\subsubsection{Ограниченный набор аналитических проходов}

Компиляторы LLVM и GCC опираются на развитые анализы: точный анализ псевдонимов указателей (aliasing), скалярную эволюцию, анализ зависимостей по данным, анализ циклов. Ранний SLP-векторизатор LLVM активно использует ScalarEvolution для проверки смежности адресов и AliasAnalysis для оценки безопасности перемещения инструкций памяти.

Инфраструктура компилятора Go, напротив, сознательно минимизирована: полноценный анализ псевдонимов отсутствует, единственным источником сведений о смежности адресов являются проходы \texttt{memcombine} и \texttt{pair}, работающие с простыми локальными шаблонами. Перенос даже ранней версии LLVM потребовал бы либо воссоздания значительной части аналитической подсистемы, либо замены точных проверок консервативными эвристиками.

\subsubsection{Явное моделирование памяти в Go SSA}

В LLVM IR порядок операций с памятью не задан жёстко; две независимые загрузки из заведомо различных областей могут быть переставлены, а \texttt{store}-инструкции не связываются виртуальными рёбрами. В Go SSA каждая операция, затрагивающая память, явно принимает и порождает значение типа \texttt{Memory}, образуя упорядоченную цепочку состояний. Эта модель, удобная для потокового анализа и понижения до машинного кода, становится серьёзным препятствием для восходящей стратегии SLP, принятой в LLVM. Две потенциально векторизуемые загрузки могут иметь различные входные значения \texttt{Memory} из-за наличия между ними записи по гарантированно не пересекающемуся адресу, что с точки зрения консервативного анализа делает их зависимыми и не подлежащими упаковке.

\subsubsection{Философия быстрой компиляции}

Компилятор Go проектировался с приоритетом скорости сборки, поэтому в нём не приветствуются дорогостоящие анализы, способные замедлить трансляцию крупных проектов. Ранний SLP-векторизатор LLVM, при всей своей относительной простоте, уже включает рекурсивный обход дерева с квадратичными проверками и нетривиальной моделью стоимости. Прямой перенос подобного прохода без адаптации к ограничениям по времени компиляции противоречил бы проектным принципам Go, требуя аккуратного выбора алгоритмических компромиссов.

\subsubsection{Итог}

Перечисленные факторы в совокупности делают невозможным простое портирование существующей реализации SLP-векторизатора. Вместо этого требуется самостоятельная разработка, опирающаяся на базовые идеи метода, но адаптированная к модели памяти, инфраструктуре анализа и ограничениям для быстрой компиляции, присущим компилятору Go.
